% USENIX Security format
\documentclass[letterpaper,twocolumn,10pt]{article}
\usepackage{usenix-2020-09}
\usepackage{graphicx}
\usepackage{booktabs}
\usepackage{amsmath,amssymb}
\usepackage{hyperref}
\usepackage{xspace}
\usepackage{listings}
\usepackage{xcolor}
\usepackage{balance}

\newcommand{\sys}{zkMetal\xspace}
\newcommand{\ie}{\textit{i.e.}\xspace}
\newcommand{\eg}{\textit{e.g.}\xspace}

\begin{document}

\title{\sys: Exploiting Unified Memory and 32-bit ALUs for\\Zero-Knowledge Cryptography on Apple Silicon GPUs}

\author{
{\rm Anonymous}\\
Under Submission
}

\maketitle

% ============================================================================
% ABSTRACT
% ============================================================================
\begin{abstract}
Apple Silicon's unified memory architecture and 32-bit GPU ALU create a
fundamentally different optimization landscape for zero-knowledge
cryptography than discrete GPUs with native 64-bit integer units. We present
\sys, the first comprehensive ZK cryptography library for Metal GPUs,
covering 60+ primitives across 18 fields and 10 elliptic curves. We report
three categories of findings: (1)~a previously undocumented M-series GPU
scheduling pathology where specific threadgroup counts cause 10--30$\times$
performance cliffs, requiring workload-aware configuration tables; (2)~a
systematic methodology for field-hardware co-design, showing that 32-bit ALU
limitations become advantages for small fields (BabyBear 8.5B~elem/s NTT,
Mersenne31 fused kernels) while requiring CIOS Montgomery with explicit
carry management for 256-bit fields; (3)~a C/NEON acceleration layer that
achieves 10--134$\times$ over baseline Swift across all CPU primitives.
Against ICICLE-Metal (the only other Metal ZK library), \sys is
30--500$\times$ faster on NTT and 6--40$\times$ faster on MSM. We
characterize performance relative to hardware theoretical floors, finding
primitives ranging from 1.2$\times$ (BabyBear NTT) to 9$\times$ (BN254 MSM)
of limits, and catalog 30+ dead-end optimizations as negative results.
\end{abstract}

% ============================================================================
% SECTION 1: INTRODUCTION
% ============================================================================
\section{Introduction}
\label{sec:intro}

Zero-knowledge proof systems have transitioned from theoretical curiosities
to deployed infrastructure. ZK rollups secure billions of dollars in
on-chain assets~\cite{zksync,starknet}, client-side proving enables
privacy-preserving identity~\cite{semaphore,anon-aadhaar}, and verifiable
computation is becoming a standard component of distributed
systems~\cite{risc0,sp1}. As these applications mature, the computational
cost of proof generation---dominated by multi-scalar multiplication (MSM),
number-theoretic transforms (NTT), and cryptographic hashing---becomes the
primary bottleneck.

The ZK acceleration landscape has focused almost exclusively on NVIDIA GPUs.
Libraries such as ICICLE~\cite{icicle}, cuZK~\cite{cuzk}, and ZPrize
submissions~\cite{zprize2022,zprize2023} target CUDA, exploiting native
64-bit integer multiply, high memory bandwidth (900+~GB/s on A100), and a
mature compiler toolchain. This focus is understandable: discrete GPUs offer
raw throughput that integrated GPUs cannot match. However, it leaves a
significant gap.

Apple Silicon now powers over 50\% of the premium laptop and tablet
market~\cite{apple-market-share}. MacBooks, iPads, and iPhones increasingly
serve as targets for client-side ZK proving---mobile wallets, browser-based
provers, and edge verification nodes. These devices expose GPU compute only
through Apple's Metal API. No mature, high-performance ZK library exists for
this platform. The sole alternative, ICICLE-Metal~\cite{icicle}, is a
port of a CUDA-first design and exhibits fundamental performance
problems: approximately 85~ms of per-call overhead (attributable to license
server communication) and limited primitive coverage.

Apple Silicon's GPU architecture differs from discrete GPUs in three ways
that materially affect ZK workloads. First, \textbf{unified memory}
eliminates PCIe transfer overhead entirely---CPU and GPU share physical
DRAM, enabling zero-copy buffer sharing. Second, the GPU provides only
\textbf{32-bit integer ALUs}: the \texttt{mulhi(ulong,ulong)} intrinsic
decomposes internally into four 32-bit multiply-add (MAD) instructions,
making 256-bit field arithmetic approximately $4\times$ more expensive per
limb operation than on CUDA hardware with native 64-bit integer multiply.
Third, the \textbf{tile-based deferred rendering} GPU architecture imposes
scheduling constraints with no public documentation, leading to performance
pathologies we characterize in Section~\ref{sec:pathology}.

These architectural differences are not merely quantitative---they create a
qualitatively different optimization landscape. Techniques that are
effective on CUDA (\eg 64-bit limb representations, aggressive kernel
fusion, reliance on high memory bandwidth to hide latency) may be
counterproductive or irrelevant on Metal. Conversely, Apple Silicon's
unified memory and 32-bit ALU create opportunities that CUDA-centric
designs miss: zero-copy buffer sharing eliminates an entire class of data
movement, and 31-bit fields such as BabyBear~\cite{plonky3} and
Mersenne31~\cite{circle-stark} map directly to single-word native
arithmetic, yielding throughput that is \emph{super-linear} relative to the
field size reduction.

\paragraph{Contributions.} We present \sys, the first comprehensive
ZK cryptography library for Apple Silicon Metal GPUs. Our contributions are:

\begin{enumerate}
\item \textbf{GPU scheduling pathology.} We document a previously unknown
  M-series GPU scheduling phenomenon where specific threadgroup counts
  (\eg powers of two and their neighbors) cause 10--30$\times$
  performance cliffs in MSM bucket accumulation. This affects all M-series
  generations and is not described in Apple's documentation
  (Section~\ref{sec:pathology}).

\item \textbf{Field-hardware co-design.} We develop a systematic
  methodology for mapping ZK field arithmetic to Metal's 32-bit ALU,
  showing that 256-bit fields require CIOS Montgomery with explicit carry
  management while 31-bit fields achieve near-hardware-floor throughput.
  BabyBear NTT at $2^{24}$ elements runs in 2.0~ms (8.5B~elem/s), within
  1.2$\times$ of the memory bandwidth limit
  (Section~\ref{sec:architecture}).

\item \textbf{Comprehensive library and evaluation.} \sys implements 60+
  primitives across 18 fields and 10 elliptic curves, including complete
  proof systems (Circle STARK, Plonk, Groth16, GKR, HyperNova, Spartan),
  post-quantum schemes (Kyber, Dilithium), and application primitives
  (ECDSA, BLS, EdDSA). Against ICICLE-Metal, \sys is 33$\times$ faster on
  MSM and 30--500$\times$ faster on NTT
  (Section~\ref{sec:optimization}--\ref{sec:evaluation}).

\item \textbf{Negative result catalog.} We document 30+ optimization
  attempts that failed, organized by failure mode (hardware mismatch,
  algorithmic dead end, compiler issue), providing guidance for future Metal
  GPU optimization efforts (Section~\ref{sec:negative}).
\end{enumerate}

Table~\ref{tab:headline} summarizes headline comparisons.
Figure~\ref{fig:overview} provides a visual comparison across the primary
primitives.

\begin{table}[t]
\centering
\small
\caption{Headline performance comparison (BN254, M3 Pro).}
\label{tab:headline}
\begin{tabular}{lrrr}
\toprule
\textbf{Primitive} & \textbf{\sys} & \textbf{ICICLE-Metal} & \textbf{Arkworks CPU} \\
\midrule
MSM $2^{18}$       & 45\,ms   & 1{,}475\,ms & 266\,ms \\
NTT $2^{20}$       & 6.1\,ms  & 194\,ms     & 7{,}300\,ms\textsuperscript{*} \\
NTT BB $2^{24}$    & 2.0\,ms  & 709\,ms     & --- \\
\bottomrule
\multicolumn{4}{l}{\footnotesize \textsuperscript{*}Vanilla Swift CPU baseline (no Arkworks NTT available).}
\end{tabular}
\end{table}

\begin{figure}[t]
\centering
% Placeholder: bar chart of zkMetal vs ICICLE-Metal vs Arkworks
% across MSM, NTT (BN254), NTT (BabyBear), Keccak Merkle, Poseidon2
\fbox{\parbox{0.9\columnwidth}{\centering\vspace{2cm}
\textbf{Figure 1:} \sys vs ICICLE-Metal vs Arkworks CPU across MSM, NTT,
and hashing primitives (log scale). \sys achieves 6--500$\times$ speedups
over ICICLE-Metal and 5--1000$\times$ over CPU baselines.
\vspace{2cm}}}
\caption{Performance comparison across core ZK primitives.}
\label{fig:overview}
\end{figure}

% ============================================================================
% SECTION 2: BACKGROUND
% ============================================================================
\section{Background}
\label{sec:background}

\subsection{Apple Silicon GPU Architecture}
\label{sec:apple-gpu}

Apple's M-series system-on-chip integrates CPU and GPU cores with a
\emph{unified memory} architecture: both processors share the same physical
DRAM, eliminating the PCIe bus that separates host and device memory on
discrete GPU systems. On the M3~Pro used for our evaluation, the shared
memory subsystem provides approximately 150~GB/s of bandwidth to both CPU
and GPU.

The M-series GPU is a tile-based deferred rendering (TBDR) architecture
originally designed for graphics workloads. For compute shaders---the
execution model used by \sys---the relevant characteristics are:

\begin{itemize}
\item \textbf{32-bit integer ALU.} The GPU provides no native 64-bit
  integer multiply. The Metal Shading Language (MSL) intrinsic
  \texttt{mulhi(ulong,ulong)} decomposes internally into four 32-bit
  multiply-add instructions. A single 256-bit Montgomery multiplication
  (CIOS~\cite{cios}) using 8$\times$32-bit limbs therefore requires
  approximately 64 MAD instructions.

\item \textbf{SIMD width 32.} Threads execute in groups of 32 (analogous
  to NVIDIA warps). SIMD shuffle operations enable intra-group
  communication but are designed for small types (4--8 bytes); shuffling
  96-byte elliptic curve points incurs substantial overhead.

\item \textbf{Threadgroup shared memory: 32\,KB.} Shared memory per
  threadgroup is limited to 32\,KB, constraining in-place NTT butterfly
  computations and requiring multi-pass algorithms for large transforms.

\item \textbf{Dispatch model.} Work is submitted through \emph{command
  buffers} containing one or more \emph{command encoders}, each encoding
  \emph{compute pipeline} dispatches. Each command buffer commit incurs
  approximately 0.5~ms of overhead, making fine-grained dispatch batching
  critical for small workloads.
\end{itemize}

\noindent
For comparison, an NVIDIA A100 provides native 64-bit integer multiply,
80\,GB of dedicated HBM2e at 2\,TB/s bandwidth, 164\,KB shared memory per
SM, and a compiler toolchain (NVCC/PTX) that exposes low-level scheduling
control. The RTX~3090~Ti, a consumer-grade comparison point, provides
approximately 900~GB/s memory bandwidth and achieves MSM~$2^{16}$ in
${\sim}$9\,ms~\cite{icicle}---roughly $3\times$ faster than \sys's
27\,ms, consistent with the $6\times$ memory bandwidth ratio modulated by
the native 64-bit multiply advantage.

\subsection{ZK Cryptographic Primitives}
\label{sec:zk-primitives}

Zero-knowledge proof systems are built from a small set of computationally
intensive primitives. We briefly describe those most relevant to our
optimization effort.

\paragraph{Multi-scalar multiplication (MSM).} Given $n$ elliptic curve
points $\{P_i\}$ and scalars $\{s_i\}$, compute $\sum_i s_i \cdot P_i$.
The Pippenger algorithm~\cite{pippenger} decomposes each scalar into
$w$-bit windows, sorts points into $2^w$ buckets per window, accumulates
each bucket, and combines via a Horner-like evaluation. MSM dominates
KZG~\cite{kzg} commitment schemes, Groth16~\cite{groth16} proving, and
IPA~\cite{bulletproofs} computations.

\paragraph{Number-theoretic transform (NTT).} The NTT is the
finite-field analogue of the FFT, computing polynomial evaluation at
roots of unity in $O(n \log n)$ field operations. NTT dominates polynomial
multiplication in STARK~\cite{stark} and Plonk~\cite{plonk} provers, FRI
commitment phases, and lattice-based cryptography.

\paragraph{Cryptographic hashing.} Algebraic hashes such as
Poseidon2~\cite{poseidon2} operate natively over finite fields, enabling
efficient Merkle tree construction within proof systems. Binary hashes
(Keccak-256~\cite{keccak}, Blake3~\cite{blake3}, SHA-256) are used for
Fiat-Shamir transcripts and data availability.

\paragraph{FRI and Sumcheck.} The FRI protocol~\cite{fri} provides
proximity testing for Reed-Solomon codes, forming the core of STARK proof
systems. The sumcheck protocol~\cite{sumcheck} reduces multivariate
polynomial evaluation claims to a sequence of univariate checks, and is
central to GKR~\cite{gkr}, Spartan~\cite{spartan}, and
HyperNova~\cite{hypernova}.

\paragraph{Montgomery arithmetic.} All field operations use Montgomery
representation~\cite{montgomery} with the CIOS (Coarsely Integrated
Operand Scanning) algorithm~\cite{cios} for modular multiplication. For a
256-bit prime, CIOS with 8$\times$32-bit limbs requires 64 multiply-add
operations with explicit carry propagation. For 31-bit primes (BabyBear,
Mersenne31), a single 32-bit MAD suffices.

\subsection{Related Work}
\label{sec:related}

\paragraph{CUDA-based ZK acceleration.} ICICLE~\cite{icicle} is the most
widely used GPU ZK library, supporting MSM, NTT, and Poseidon hashing on
NVIDIA GPUs. cuZK~\cite{cuzk} optimizes MSM via an elliptic curve-specific
bucket sorting approach. ZPrize~\cite{zprize2022,zprize2023} competitions
have driven MSM optimization on both FPGA and GPU targets. These systems
target CUDA and exploit native 64-bit integer arithmetic, high memory
bandwidth, and mature profiling tools (Nsight, PTX inspection).

\paragraph{Metal-based ZK.} ICICLE-Metal~\cite{icicle} ports the ICICLE
library to Apple's Metal API. In our measurements (v3.8.0), ICICLE-Metal
exhibits approximately 85~ms per-call overhead (attributable to license
server initialization) and covers only MSM and NTT. At BN254 MSM~$2^{18}$,
ICICLE-Metal achieves 1{,}475~ms versus \sys's 45\,ms---a 33$\times$ gap.
MoPro~v2~\cite{mopro} reports MSM benchmarks on M3~Air (678~ms at
$2^{18}$), using a different methodology.

\paragraph{CPU-based libraries.} Arkworks~\cite{arkworks} is the standard
Rust-based CPU ZK library, achieving 266~ms for BN254 MSM at $2^{18}$ and
69~ms at $2^{16}$ on M3-class hardware. gnark~\cite{gnark} (Go) and
Plonky3~\cite{plonky3} (Rust, BabyBear-focused) provide additional CPU
baselines. Our C~CIOS CPU path achieves 68~ms for BN254 MSM at $2^{16}$,
matching Arkworks' 69~ms at the same size, demonstrating that systematic
field arithmetic optimization alone closes the gap with state-of-the-art
CPU implementations.

\paragraph{Gap.} No prior work systematically characterizes the Apple
Silicon GPU for ZK workloads, documents its scheduling pathologies, or
develops a field-hardware co-design methodology for its 32-bit ALU. \sys
fills this gap with the most comprehensive Metal ZK implementation by an
order of magnitude (60+ primitives versus ICICLE-Metal's 2).

% ============================================================================
% SECTION 3: ARCHITECTURE
% ============================================================================
\section{System Architecture}
\label{sec:architecture}

\sys is organized as a three-tier execution model that routes computation to
the most efficient execution unit based on workload size and characteristics.
Figure~\ref{fig:architecture} illustrates the architecture and data flow.

\begin{figure}[t]
\centering
% Placeholder: architecture diagram showing three tiers
% Tier 1 (Metal GPU), Tier 2 (C/NEON CPU), Tier 3 (Swift orchestration)
% Annotate unified memory zero-copy paths between tiers
\fbox{\parbox{0.9\columnwidth}{\centering\vspace{2.5cm}
\textbf{Figure 2:} Three-tier execution model. Metal GPU shaders (Tier~1)
handle compute-intensive parallel workloads. C/NEON kernels (Tier~2)
provide optimized CPU paths. Swift orchestration (Tier~3) manages pipeline
selection and buffer lifecycle. Unified memory enables zero-copy data sharing
between all tiers.
\vspace{2.5cm}}}
\caption{Three-tier \sys architecture with unified memory data flow.}
\label{fig:architecture}
\end{figure}

\subsection{Three-Tier Execution Model}
\label{sec:tiers}

\paragraph{Tier 1: Metal GPU shaders.} Compute-intensive parallel
kernels---MSM bucket accumulation, NTT butterfly operations, hash batch
evaluation, polynomial operations---are implemented as Metal compute
shaders written in Metal Shading Language (a C++14 dialect). Each
\emph{engine} class (\eg \texttt{MetalMSM}, \texttt{NTTEngine},
\texttt{Keccak256Engine}) encapsulates pipeline state objects, manages
buffer allocation, and encodes dispatches into command buffers.

Buffer management follows a \emph{grow-only caching} pattern: each engine
maintains a cache of Metal buffers indexed by role. On first invocation, the
engine allocates buffers sized to the workload. On subsequent invocations,
buffers are reused if they are large enough, or reallocated at the new
(larger) size. Buffers are never shrunk. This eliminates repeated allocation
overhead---a significant cost on Metal, where buffer creation involves
kernel-level page mapping. Caching alone improves FRI engine throughput by
2.4$\times$ and Keccak engine throughput by 28\%.

Command buffer batching is equally critical. Each Metal command buffer
commit incurs approximately 0.5~ms of fixed overhead. For multi-pass
algorithms (\eg NTT with $\log_2 n$ butterfly stages, FRI with multiple
fold rounds), encoding all dispatches into a single command buffer via a
single compute command encoder amortizes this overhead. For BN254 NTT at
$2^{16}$, the single-encoder pattern yields a $2\times$ improvement over
per-pass command buffer submission.

\paragraph{Tier 2: C/NEON CPU kernels.} When GPU dispatch overhead
exceeds the compute savings---typically for inputs smaller than
$2^{12}$--$2^{14}$ elements---\sys routes computation to optimized C
functions linked via Swift Package Manager C targets. These functions use
\texttt{\_\_uint128\_t} for 64-bit CIOS Montgomery multiplication on ARM64
(exploiting the processor's native \texttt{MUL}/\texttt{UMULH} instruction
pair) and NEON SIMD for small-field operations (4-wide \texttt{uint32}
Montgomery butterfly for BabyBear).

The C~CIOS layer delivers 10--134$\times$ speedup over equivalent Swift
implementations across all CPU primitives. The gap is not primarily due to
Swift's optimizer quality for field arithmetic---Swift achieves 80--90\% of
C performance for isolated 256-bit multiply---but rather to Swift's
function call overhead in tight inner loops (Pippenger bucket accumulation,
NTT butterfly iteration), where the overhead per scalar operation dominates
the actual computation.

\paragraph{Tier 3: Swift orchestration.} The top-level engine layer,
written in Swift, manages execution path selection (GPU vs.\ CPU based on
input size), auto-tuning of GPU parameters (threadgroup sizes, MSM window
widths, NTT split thresholds), buffer lifecycle, and command buffer
construction. Auto-tuning results are persisted to disk and reused across
invocations, re-triggering calibration only when the GPU hardware changes.

\subsection{Field Arithmetic on 32-bit ALU}
\label{sec:field-arith}

The Metal GPU's lack of native 64-bit integer multiply creates a
field-size-dependent performance landscape that fundamentally shapes our
optimization strategy.

\paragraph{256-bit fields (BN254, BLS12-381).} Field elements are
represented as 8$\times$32-bit limbs in Montgomery form. Each
\texttt{fp\_mul} invocation executes the CIOS algorithm with explicit carry
chains: approximately 64 MAD instructions, with carries propagated after
each inner-loop iteration to prevent overflow of the 32-bit accumulators.
Sum-of-Squares (SOS) squaring optimization exploits the symmetry of
$a \times a$ to reduce the multiply count by approximately 25\%, saving
24~ms in proof-level benchmarks (101$\to$77~ms in the Barretenberg
UltraHonk context).

\paragraph{64-bit fields (Goldilocks).} The Goldilocks prime
$p = 2^{64} - 2^{32} + 1$ admits a special reduction: overflow beyond
$2^{64}$ can be handled by a single-word correction, avoiding the full
CIOS carry chain. On Metal's 32-bit ALU, this still requires two 32-bit
limbs, but the reduction step is substantially cheaper than for general
256-bit primes.

\paragraph{31-bit fields (BabyBear, Mersenne31).} These fields fit in a
single 32-bit word. Montgomery multiplication reduces to a single MAD
instruction---natively supported by the Metal ALU at full throughput. This
is where Metal's 32-bit ALU transforms from a limitation into an
\emph{advantage}: a 64-bit CUDA GPU wastes half its ALU width on
single-word multiplies, while Metal utilizes its full datapath. The
density advantage is super-linear: BabyBear NTT at $2^{24}$ achieves
8.5B~elem/s versus BN254's 144M~elem/s---a 59$\times$ throughput ratio
for an 8$\times$ element size reduction, reflecting both the reduced
per-element compute and improved memory access patterns (4-byte elements
versus 32-byte elements per cache line).

\paragraph{Binary tower fields (Binius).} Addition in binary extension
fields is XOR, which is essentially free on GPU hardware. Multiplication
uses lookup tables. Metal's 32-bit ALU is well suited to these
table-driven operations, and the binary tower
$\text{GF}(2^8) \to \text{GF}(2^{128})$ maps naturally to
packed 32-bit word operations.

\subsection{Unified Memory Exploitation}
\label{sec:unified-mem}

Apple Silicon's unified memory is the single most consequential
architectural difference from discrete GPU platforms. CPU and GPU access the
same physical DRAM pages without any explicit transfer. \sys exploits this
in three ways.

\paragraph{Zero-copy buffer sharing.} When the CPU prepares MSM scalars
and elliptic curve points, the GPU reads them directly from the same memory
without a copy or DMA transfer. This eliminates the PCIe transfer overhead
that dominates small-batch GPU workloads on discrete platforms. For MSM at
$2^{14}$ points, where the input data is approximately 4\,MB, a PCIe~4.0
transfer would add ${\sim}$0.2~ms---comparable to the compute time itself.
On \sys, this cost is zero.

\paragraph{Grow-only buffer caching.} As described in
Section~\ref{sec:tiers}, all engines cache Metal buffers across invocations.
Because unified memory means buffer contents are visible to both CPU and
GPU without marshaling, the caching strategy is simpler than on discrete
platforms: no coherence management, no pinned memory allocation, no
staging buffers. The measured impact is substantial: FRI engine throughput
improves by 2.4$\times$ and Keccak engine throughput by 28\% from caching
alone.

\paragraph{Single command buffer batching.} Multi-pass GPU algorithms
(\eg NTT butterfly stages, FRI fold rounds, MSM multi-window accumulation)
encode all dispatches into a single command buffer. This amortizes the
${\sim}$0.5~ms per-command-buffer overhead that would otherwise dominate
multi-pass algorithms with small per-pass compute. For NTT at $2^{16}$ with
16 butterfly stages, this overhead would otherwise constitute 8~ms of
pure scheduling cost---exceeding the 0.47~ms total compute time by
$17\times$.

The combination of zero-copy sharing, buffer caching, and command buffer
batching eliminates three classes of overhead (transfer, allocation,
scheduling) that are inherent costs on discrete GPU platforms. This
partially compensates for Apple Silicon's lower raw compute throughput
(3.6~TFLOPS versus 40+~TFLOPS on an A100) and lower memory bandwidth
(150~GB/s versus 2~TB/s). For workloads where these overheads dominate
(\eg small-to-medium MSM, iterative protocols with many round-trips), \sys
can be competitive with or faster than discrete GPU implementations
despite the raw hardware gap.

\end{document}
